\section{\textit{Decision Tree Learning} (C4.5)}
\label{sec:dtl}

\textit{Decision tree learning} (DTL) membangun pohon aturan ``jika-maka'' dengan memecah data secara
rekursif pada atribut yang paling memisahkan kelas. Varian yang kami implementasikan adalah C4.5
\cite{quinlan1993c45}: atribut dipilih dengan \textit{gain ratio}, atribut kontinu dipecah dengan
ambang, atribut kategorikal dipecah satu cabang per kategori, dan nilai hilang ditangani dengan bobot
pecahan. Karena pemecahannya hanya bergantung pada urutan nilai, pohon tidak memerlukan penskalaan
fitur, tetapi mudah \textit{overfitting} jika dibiarkan tumbuh sampai murni. Implementasi kami
menyediakan \textit{reduced error pruning} dari materi kelas \cite{slide-dtl-issues} serta
\textit{pessimistic pruning} \cite{quinlan1993c45} dan \textit{cost-complexity pruning}
\cite{breiman1984cart} sebagai bonus algoritma optimasi.

Konfigurasi terbaik hasil \textit{stratified} 5-\textit{fold cross validation} adalah pohon C4.5
dengan kedalaman maksimum 20, minimal 2 contoh per cabang, \textit{cost-complexity pruning}
$\alpha = 2{,}5$, 46 fitur, dan \textit{threshold} 0,4. Konfigurasi ini mencapai Macro F1
$0{,}8692 \pm 0{,}0030$ (rata-rata per \textit{fold}, \textit{seed} 42) dan 0,87393 pada
\textit{public leaderboard} Kaggle. Faktor penentunya adalah \textit{pruning}: pohon yang sama tanpa
\textit{pruning} memiliki sekitar 688 daun dan hanya mencapai 0,8497, sedangkan dengan
\textit{pruning} pohon tinggal sekitar 66 daun.

\subsection{Implementasi \textit{From Scratch}}
\label{subsec:dtl-implementasi}

\subsubsection{Landasan Teori}
\todo{Jelaskan entropi, \textit{information gain}, \textit{split information}, \textit{gain ratio}; pencarian \textit{threshold} untuk atribut kontinu; penanganan \textit{missing value} dengan bobot fraksional; dan \textit{pruning} (\textit{pessimistic error pruning} C4.5). Acu \cite{slide-dtl,slide-dtl-issues,quinlan1993c45}. Jelaskan juga alasan memilih C4.5 dibanding ID3 dan CART. Sertakan rumus utama dalam lingkungan \texttt{equation}/\texttt{align}.}

\subsubsection{Struktur Kelas dan Fungsi}
\todo{Deskripsikan kelas pada \texttt{src/algoritma/dtl/}: tiap modul (\texttt{criteria.py}, \texttt{splitting.py}, \texttt{tree.py}, \texttt{pruning.py}, \texttt{serialization.py}, \texttt{dtl.py}) beserta fungsi pentingnya dan parameter kelas \texttt{DTL}. Isi tabel di bawah, satu baris per fungsi/metode.}

\begin{table}[H]
	\centering
	\caption{Deskripsi kelas dan fungsi DTL}
	\label{tab:dtl-fungsi}
	\begin{tabularx}{\textwidth}{l Y}
		\toprule
		\textbf{Fungsi/Metode} & \textbf{Deskripsi} \\
		\midrule
		\texttt{nama\_fungsi} & \todo{deskripsi} \\
		\bottomrule
	\end{tabularx}
\end{table}

\subsubsection{\textit{Source Code}}
\todo{Sisipkan potongan kode inti (perhitungan \textit{gain ratio}, pembangunan pohon rekursif, dan \textit{pruning}) dengan makro \texttt{\textbackslash kode}, lalu jelaskan alurnya. Cek ulang nomor baris jika kode berubah.}
% \kode{../../src/algoritma/dtl/dtl.py}{1}{30}{Judul potongan kode}{lst:dtl-inti}

\subsubsection{\textit{Save} dan \textit{Load} Model}
\todo{Format penyimpanan (JSON/\texttt{.npz}/\textit{pickle}), isi yang disimpan, dan contoh pemakaian.}

\subsection{\textit{Data Cleaning} dan \textit{Preprocessing}}
\label{subsec:dtl-preprocessing}

Seluruh langkah diimplementasikan pada \nolinkurl{src/kaggle/booking_features.py} (pandas dan NumPy)
dan dipanggil oleh notebook \nolinkurl{src/kaggle/dtl_notebook.ipynb} dengan urutan berikut:
\begin{enumerate}
	\item pembuangan kolom \texttt{id} dan duplikat (data latih saja, \texttt{clean}),
	\item fitur tanggal pemesanan (\texttt{add\_booking\_dates}),
	\item imputasi dan penanda kekosongan (\texttt{impute\_and\_flag}),
	\item fitur grup pemesanan (\texttt{sibling\_pairs} dan \texttt{build\_features}),
	\item fitur lama menginap,
	\item \textit{one-hot encoding}.
\end{enumerate}
Semua statistik (median, modus, daftar kategori, dan label untuk fitur grup) dipelajari dari data
latih saja, lalu diterapkan ke data validasi dan data uji \cite{slide-dataprep}. Fungsi
\texttt{build\_features} menerima data latih dan data lain sebagai dua argumen terpisah, sehingga
kedua data tidak pernah digabung. Pada validasi silang, fungsi ini dipanggil ulang di setiap
\textit{fold} dengan bagian latih \textit{fold} tersebut sebagai data latih. Angka analisis di bawah
berasal dari seluruh data latih bersih (27.309 baris).

\subsubsection{Penanganan \textit{Missing Values}}

Delapan fitur memiliki nilai kosong, dan hanya 60,9\% baris data latih yang lengkap.
Tabel~\ref{tab:dtl-missing} membandingkan tingkat pembatalan saat kolom kosong dan saat terisi.

\begin{table}[H]
	\centering
	\caption{Kolom dengan nilai kosong pada data latih bersih dan data uji, serta tingkat pembatalan saat kolom kosong dan terisi}
	\label{tab:dtl-missing}
	\small
	\begin{tabular}{l r c c c c}
		\toprule
		\textbf{Kolom} & \textbf{Kosong} & \textbf{\% latih} & \textbf{\% uji} & \textbf{Batal (kosong)} & \textbf{Batal (terisi)} \\
		\midrule
		\texttt{num\_adults} & 809 & 3,0 & 2,6 & 0,335 & 0,327 \\
		\texttt{num\_children} & 3.035 & 11,1 & 10,9 & 0,328 & 0,327 \\
		\texttt{meal\_package} & 1.760 & 6,4 & 6,5 & 0,319 & 0,328 \\
		\texttt{needs\_parking} & 1.082 & 4,0 & 3,6 & 0,312 & 0,328 \\
		\texttt{room\_category} & 1.002 & 3,7 & 3,5 & 0,348 & 0,327 \\
		\texttt{days\_before\_arrival} & 1.120 & 4,1 & 4,2 & 0,312 & 0,328 \\
		\texttt{nightly\_rate\_idr} & 1.992 & 7,3 & 7,9 & 0,309 & 0,329 \\
		\texttt{special\_requests} & 2.499 & 9,2 & 8,7 & 0,409 & 0,319 \\
		\bottomrule
	\end{tabular}
\end{table}

Proporsi kosong di data latih dan data uji hampir sama. Pada tujuh kolom, tingkat pembatalan saat
kosong hanya berbeda paling banyak 0,021 dari saat terisi. Pengecualiannya
\texttt{special\_requests}: baris yang kosong batal 0,409, jauh di atas baris terisi (0,319).

Kelas \texttt{DTL} dapat menerima nilai kosong secara langsung dengan bobot pecahan
(Subbab~\ref{subsec:dtl-implementasi}). Meski begitu, \textit{pipeline} final mengimputasi nilai
kosong dengan dua alasan. Pertama, varian yang menerima data mentah jauh lebih lambat karena baris
bernilai kosong diteruskan ke semua cabang, sehingga hanya dapat diuji sampai kedalaman 10
(Subbab~\ref{subsec:dtl-validasi}). Kedua, data yang sudah lengkap dapat dipakai langsung oleh
pembanding \textit{scikit-learn} dan PCA, yang tidak menerima nilai kosong, sehingga semua
perbandingan memakai matriks fitur yang sama. Langkahnya sebagai berikut.
\begin{itemize}
	\item \textbf{Imputasi.} Fitur numerik diisi median dan fitur kategorikal diisi modus data latih.
	\item \textbf{Penanda kekosongan.} Setiap kolom yang memiliki nilai kosong di data latih mendapat
	kolom biner \texttt{missing\_<kolom>}, termasuk fitur turunan \texttt{book\_ord} yang ikut kosong
	ketika \texttt{days\_before\_arrival} kosong. Hasilnya sembilan penanda. Penanda ini menjaga
	informasi yang hilang akibat imputasi, terutama pada \texttt{special\_requests}.
\end{itemize}

\subsubsection{Data Duplikat dan Data Tidak Valid}

\paragraph{Duplikat.} Kolom \texttt{id} dibuang karena hanya berupa \textit{hash} acak. Kolom ini
dibaca sebagai teks karena sebagian nilainya, misalnya \texttt{123e4567}, akan terbaca sebagai
notasi ilmiah jika dibaca sebagai angka. Data latih memuat 64 baris duplikat persis (seluruh fitur
dan label sama), sehingga tersisa 27.309 baris. Ada juga dua pasang baris dengan fitur identik
tetapi label berbeda. Keempat baris dipertahankan karena tidak ada dasar untuk memilih salah satu
label. Data uji tidak diubah barisnya karena setiap baris harus diprediksi.

\paragraph{Data tidak valid.} Terdapat 106 baris dengan \texttt{num\_adults} = 0 dan 408 baris
dengan \texttt{nightly\_rate\_idr} = 0. Kedua nilai ini mungkin terjadi pada pemesanan nyata, misalnya
kamar komplimen, sehingga dianggap nilai asli dan dipertahankan.

\paragraph{\textit{Outlier}.} Tidak ada nilai yang dipangkas. Pemecahan pohon hanya bergantung pada
urutan nilai terhadap ambang, sehingga nilai ekstrem tidak menarik ambang seperti pada model linear.

\subsubsection{Konstruksi Fitur dan \textit{Encoding}}

\paragraph{Fitur tanggal pemesanan.} \texttt{arr\_ord} adalah tanggal kedatangan dalam jumlah hari
sejak 1 Januari 2023, \texttt{book\_ord} = \texttt{arr\_ord} $-$ \texttt{days\_before\_arrival}
adalah perkiraan tanggal pemesanan, dan \texttt{arr\_dow} adalah hari kedatangan dalam pekan.
Kolom tahun, bulan, dan tanggal yang terpisah tidak dapat dibandingkan dalam satu pemecahan, sedangkan
\texttt{arr\_ord} menyatakan urutan waktu dalam satu angka.

\paragraph{Fitur grup pemesanan.} Sebagian pemesanan memiliki pemesanan lain di data latih dengan
tanggal kedatangan, lama menginap, dan kanal yang sama, yang kemungkinan berasal dari satu rombongan.
Pohon biasa tidak dapat memakai pola ini karena setiap baris diproses sendiri-sendiri. Untuk setiap
baris, dicari baris data latih lain (bukan baris itu sendiri) yang kemungkinan satu grup dengan aturan
berikut:
\begin{itemize}
	\item tanggal kedatangan, jumlah malam akhir pekan, jumlah malam hari kerja, dan kanal pemesanan
	sama persis;
	\item selisih \texttt{days\_before\_arrival} paling banyak 10 hari;
	\item selisih harga per malam paling banyak 5\% dari harga yang lebih tinggi di antara keduanya;
	\item tipe kamar, paket makan, jumlah tamu, \textit{special request}, dan kebutuhan parkir sama,
	atau kosong di salah satu baris.
\end{itemize}
Baris segrup $j$ diberi bobot $w_j = \max(0,\ 1 - |\Delta_j| / 11)$, dengan $\Delta_j$ selisih
\texttt{days\_before\_arrival}, atau $w_j = 0{,}1$ jika salah satu nilainya kosong. Dari bobot ini
dibentuk dua fitur. \texttt{sib\_n} adalah jumlah bobot baris segrup, dan \texttt{sib\_te} adalah
rata-rata label berbobot yang dihaluskan menuju rata-rata label global (\textit{smoothed target
encoding}) \cite{micci2001targetenc}:
\begin{equation}
	\texttt{sib\_te} = \frac{\sum_{j \in G} w_j\, y_j + m \cdot \bar{y}}{\sum_{j \in G} w_j + m}, \qquad m = 1
	\label{eq:dtl-sibling}
\end{equation}
dengan $G$ himpunan baris segrup dan $\bar{y}$ proporsi kelas batal di data latih.

Fitur yang memakai label rawan kebocoran. Karena itu, baris data latih dibagi secara acak ke dalam
lima \textit{fold} internal, dan \texttt{sib\_te} untuk sebuah baris hanya memakai label baris segrup
dari \textit{fold} internal lain. Data uji memakai label seluruh data latih, dan \texttt{sib\_n} untuk
data uji dikali $4/5$ karena acuannya seluruh data latih, bukan empat dari lima \textit{fold}.
Hasilnya, 39,3\% baris data latih dan 42,1\% baris data uji memiliki minimal satu baris segrup.
Tabel~\ref{tab:dtl-sibling} menunjukkan bahwa \texttt{sib\_te} memisahkan kelas dengan jelas, dan
fitur ini menjadi akar pohon final (Subbab~\ref{subsec:dtl-bonus}).

\begin{table}[H]
	\centering
	\caption{Proporsi label batal pada data latih menurut nilai \texttt{sib\_te} (\textit{out-of-fold})}
	\label{tab:dtl-sibling}
	\begin{tabular}{c r r}
		\toprule
		\textbf{Rentang \texttt{sib\_te}} & \textbf{Jumlah baris} & \textbf{Proporsi batal} \\
		\midrule
		$(0{,}0;\ 0{,}2]$ & 4.547 & 0,026 \\
		$(0{,}2;\ 0{,}4]$ & 18.639 & 0,268 \\
		$(0{,}4;\ 0{,}6]$ & 715 & 0,733 \\
		$(0{,}6;\ 0{,}8]$ & 1.216 & 0,921 \\
		$(0{,}8;\ 1{,}0]$ & 2.192 & 0,998 \\
		\bottomrule
	\end{tabular}
\end{table}

\paragraph{Fitur lama menginap.} \texttt{nights} adalah jumlah malam akhir pekan dan malam hari kerja,
dan \texttt{weekend\_share} adalah proporsi malam akhir pekan terhadap \texttt{nights}. Pohon dengan
pemecahan sejajar sumbu tidak dapat membentuk jumlah atau rasio dua fitur dalam satu pemecahan,
sehingga kedua fitur ini disediakan langsung.

\paragraph{\textit{Encoding}.} Tiga fitur kategorikal (\texttt{meal\_package},
\texttt{room\_category}, dan \texttt{booking\_channel}) di-\textit{one-hot}, sehingga setiap kategori
menjadi kolom biner dan pohon memecahnya dengan ambang 0,5. Kelas \texttt{DTL} sebenarnya dapat
memecah kategori mentah satu cabang per kategori. \textit{One-hot} dipakai agar matriks fitur sama
untuk \texttt{DTL}, \textit{scikit-learn}, dan PCA. Kolom data uji disamakan dengan kolom data latih,
dan kategori yang tidak ada di data latih menjadi nol di semua kolom \textit{one-hot}-nya.

Setelah semua langkah di atas, model menerima 46 fitur numerik: 14 fitur numerik asli, 3 fitur
tanggal, 9 penanda kekosongan, 2 fitur grup pemesanan, 2 fitur lama menginap, dan 16 kolom
\textit{one-hot}.

\subsubsection{Normalisasi dan Standarisasi}

Fitur tidak dinormalisasi. Pemecahan pohon hanya bergantung pada urutan nilai terhadap ambang,
sehingga transformasi monoton seperti standarisasi atau \textit{min-max} tidak mengubah pohon yang
terbentuk. Standarisasi hanya dipakai sebelum PCA pada Subbab~\ref{subsec:dtl-dimred}, karena PCA
peka terhadap skala.

\subsubsection{Penanganan \textit{Imbalanced Data}}

Kelas batal mencakup 32,75\% data latih (8.944 dari 27.309 baris). Distribusi ini tidak seimbang,
tetapi tidak ekstrem. Tidak dilakukan \textit{resampling}, karena efek utamanya pada pohon adalah
menggeser proporsi kelas di setiap daun, dan pergeseran yang sama dapat diperoleh langsung dengan
menggeser \textit{threshold}. Ketidakseimbangan ditangani di sisi keputusan: evaluasi memakai Macro F1
sehingga kelas batal berbobot sama dengan kelas mayoritas, dan \textit{threshold} 0,4 dipakai alih-alih
0,5 (Subbab~\ref{subsec:dtl-validasi}). Pembobotan kelas (\texttt{class\_weight}) tersedia di
implementasi, tetapi hanya dipakai untuk \textit{cost-sensitive classification}
(Subbab~\ref{sec:cost-sensitive}), karena di sana tujuannya biaya, bukan Macro F1.

\subsubsection{\textit{Feature Selection}}
\label{subsec:dtl-fs}

Setiap fitur diberi skor berupa \textit{information gain} terbaiknya (dengan ambang terbaik) pada
bagian latih setiap \textit{fold}, lalu fitur diurutkan menurut skor tersebut. Lima fitur dengan rata-rata
\textit{gain} tertinggi adalah \texttt{sib\_te} (0,199), \texttt{days\_before\_arrival} (0,113),
\texttt{special\_requests} (0,043), \texttt{arr\_ord} (0,041), dan \texttt{arrival\_year} (0,026).
Peringkat dihitung ulang di setiap \textit{fold}, sehingga label \textit{fold} validasi tidak ikut
menentukan fitur yang dipilih. Tabel~\ref{tab:dtl-fs} menunjukkan hasilnya.

\begin{table}[H]
	\centering
	\caption{Pengaruh \textit{feature selection} dan PCA pada model final (5-\textit{fold}, \textit{seed} 42, \textit{threshold} 0,4)}
	\label{tab:dtl-fs}
	\begin{tabular}{l c c c}
		\toprule
		\textbf{Himpunan fitur} & \textbf{Macro F1} & \textbf{Rata-rata daun} & \textbf{Waktu (s)} \\
		\midrule
		Semua fitur (46) & $0{,}8692 \pm 0{,}0030$ & 65,6 & 22,2 \\
		30 fitur teratas & $0{,}8697 \pm 0{,}0036$ & 64,2 & 15,3 \\
		20 fitur teratas & $0{,}8677 \pm 0{,}0034$ & 72,2 & 12,1 \\
		12 fitur teratas & $0{,}8426 \pm 0{,}0029$ & 103,0 & 15,4 \\
		PCA 95\% variansi & $0{,}7804 \pm 0{,}0034$ & 304,8 & 89,1 \\
		\bottomrule
	\end{tabular}
\end{table}

Membuang fitur tidak memberi kenaikan yang berarti. Tiga puluh fitur teratas hanya naik 0,0005, jauh
di bawah simpangan baku antar-\textit{fold} (0,003), sedangkan 12 fitur teratas turun 0,027. Pohon
sudah memilih fitur sendiri di setiap simpul, sehingga fitur yang lemah sendirian tetap dapat dipakai
di simpul yang lebih dalam ketika berinteraksi dengan fitur lain. Model final memakai semua 46 fitur.

\subsubsection{\textit{Dimensionality Reduction}}
\label{subsec:dtl-dimred}

PCA diuji dengan mempertahankan komponen yang mencakup 95\% variansi. Sebelum PCA, fitur
distandarisasi, dan tiga fitur yang sangat miring ditransformasi dahulu (akar kuadrat untuk
\texttt{days\_before\_arrival}, logit untuk \texttt{sib\_te}, dan $\log(1+x)$ untuk \texttt{sib\_n}).
Hasilnya pada Tabel~\ref{tab:dtl-fs}: Macro F1 turun dari 0,8692 menjadi 0,7804, pohon membesar dari
sekitar 66 menjadi 305 daun, dan waktunya empat kali lebih lama.

Penurunan ini sesuai dengan cara kerja pohon. Pemecahan pohon sejajar sumbu fitur, sedangkan setiap
komponen PCA adalah kombinasi linear semua fitur. Aturan sederhana seperti ``\texttt{sib\_te}
$\le 0{,}47$'' tidak lagi dapat dinyatakan dengan satu pemecahan, sehingga pohon memerlukan banyak
pemecahan bertingkat untuk mendekatinya. Selain itu, arti setiap simpul hilang. PCA tidak dipakai pada
model final.

\subsection{\textit{Modeling} dan Validasi}
\label{subsec:dtl-validasi}

\subsubsection{Skema Validasi}
\todo{Skema validasi untuk DTL (\textit{stratified k-fold} atau \textit{train-test split}), nilai $k$/rasio, dan \textit{seed}. Jika sama dengan kerangka umum (Subbab~\ref{sec:kerangka-umum}), cukup rujuk dan jelaskan perbedaannya saja, misalnya \textit{subsampling} karena waktu komputasi.}

\subsubsection{\textit{Hyperparameter Tuning}}
\todo{Ruang pencarian (\texttt{max\_depth}, \texttt{min\_samples\_leaf}, kriteria \textit{split}, strategi \textit{pruning}, \texttt{confidence}, \texttt{ccp\_alpha}), metode pencarian, dan hasilnya. Sertakan tabel atau grafik skor CV terhadap hiperparameter.}

\begin{table}[H]
	\centering
	\caption{Hasil \textit{tuning} hiperparameter DTL}
	\label{tab:dtl-tuning}
	\begin{tabular}{l l c}
		\toprule
		\textbf{Hiperparameter} & \textbf{Ruang pencarian} & \textbf{Terbaik} \\
		\midrule
		\todo{} & & \\
		\bottomrule
	\end{tabular}
\end{table}

\subsubsection{Model Final}
\todo{Konfigurasi final DTL, skor CV (rata-rata $\pm$ simpangan baku), dan lokasi model tersimpan di \texttt{output/}.}

\subsection{Hasil dan Perbandingan dengan \textit{Scikit-learn}}
\label{subsec:dtl-hasil}

\todo{Bandingkan dengan \texttt{DecisionTreeClassifier} dalam tiga konfigurasi wajib dari spesifikasi: ID3 (\texttt{criterion='entropy'}), C4.5 (\texttt{criterion='entropy'} + \texttt{ccp\_alpha}), dan CART (\texttt{criterion='gini'}). Jelaskan \textit{preprocessing} tambahan untuk \textit{sklearn} (\textit{encoding} dan imputasi) serta cara memilih \texttt{ccp\_alpha}. Pakai \textit{preprocessing} dan \textit{fold} yang sama agar perbedaan hanya berasal dari implementasi \cite{pedregosa2011sklearn}.}

\begin{table}[H]
	\centering
	\caption{Perbandingan DTL \textit{from scratch} dan \textit{scikit-learn}}
	\label{tab:dtl-hasil}
	\begin{tabular}{l c c c c c}
		\toprule
		\textbf{Implementasi} & \textbf{Macro F1} & \textbf{Accuracy} & \textbf{Precision} & \textbf{Recall} & \textbf{Waktu (s)} \\
		\midrule
		\textit{From scratch} (C4.5) & & & & & \\
		\textit{Sklearn} ID3 & & & & & \\
		\textit{Sklearn} C4.5 (\texttt{ccp\_alpha}) & & & & & \\
		\textit{Sklearn} CART & & & & & \\
		\bottomrule
	\end{tabular}
\end{table}

\begin{figure}[H]
	\centering
	\placeholdergambar{\textit{Confusion matrix} \textit{from scratch} vs \textit{scikit-learn}}
	\caption{\textit{Confusion matrix} DTL}
	\label{fig:dtl-cm}
\end{figure}

\subsubsection{Pengaruh \textit{Pruning}}
\todo{Bandingkan strategi \textit{pruning} (tanpa, \textit{pessimistic}, \textit{reduced error}, \textit{cost-complexity}): kedalaman, jumlah daun, dan \textit{macro F1}. Sertakan grafik.}

\subsubsection{Analisis}
\todo{Bandingkan kedalaman, jumlah daun, dan kinerja. Jelaskan perbedaan karena \textit{gain ratio} vs entropi/gini, \textit{multiway} vs \textit{binary split}, penanganan \textit{null}, dan jenis \textit{pruning}.}

\subsection{Bonus DTL}
\label{subsec:dtl-bonus}

\todo{Hapus subbab bonus yang tidak dikerjakan.}

\subsubsection{Visualisasi Percabangan Pohon}
\todo{Gambar percabangan pohon beserta parameter \textit{top-N} percabangan yang ditampilkan. Jelaskan fungsi visualisasi, cara memilih \textit{top-N}, dan tampilkan gambarnya.}

\begin{figure}[H]
	\centering
	\placeholdergambar{Visualisasi Percabangan Pohon}
	\caption{Visualisasi Percabangan Pohon DTL}
	\label{fig:dtl-bonus-0}
\end{figure}

\subsubsection{Algoritma Optimasi di Luar Materi Kelas}
\todo{Jika ada: jelaskan algoritmanya, mengapa dapat mengoptimasi model, bukti perbaikannya (sebelum vs sesudah), dan \textbf{referensinya (tanpa referensi bonus dianulir)}. Misalnya \textit{cost-complexity pruning}, penalti MDL pada \textit{threshold} kontinu, \textit{shrinkage} probabilitas daun, atau Laplace \textit{correction}.}

